\subsection{Backend}
\label{subsec:kelvin-backend}

Backend systému je implementován v programovacím jazyce Python s využitím webového frameworku Django.
Django zajišťuje správu databázových modelů prostřednictvím ORM (Object-Relational Mapping), autentizaci a autorizaci uživatelů, správu souborů (odevzdaných zdrojových kódů a testovacích dat) a základní routování požadavků.

Použitým databázovým systémem je PostgreSQL, přičemž veškerý přístup k němu probíhá výhradně přes Django ORM\@.
Soubory zdrojových kódů jsou ukládány na souborový uložišti serveru, přičemž v databázi jsou uloženy pouze metadata (cesty, časové razítka, vazby na studenty a úlohy).

\subsubsection*{Evaluační pipeline}
\label{subsubsec:kelvin-evaluation-pipeline}

Stávající evaluační pipeline systému zajišťuje kompilaci odevzdaného zdrojového kódu a spuštění automatických testů v izolovaném prostředí.
Izolace je realizována prostřednictvím kontejnerů (Docker) a omezení přístupu k systémovým zdrojům, což chrání infrastrukturu před škodlivým nebo nekorektním kódem a současně garantuje jednotné podmínky evaluace pro všechny studenty.

Vyhodnocení je implementováno jako asynchronní úloha z pohledu aplikační architektury, nicméně z hlediska uživatelského chování se jedná o proces synchronní, protože student čeká na jeho dokončení před zobrazením výsledků testů.

Zařazení dalších výpočetně náročných operací přímo do této pipeline, například generování analýzy pomocí LLM, by vedlo k prodloužení doby čekání na výsledek.
I pokud by byly tyto operace realizovány technicky asynchronně, jejich začlenění do stejného evaluačního procesu by mělo přímý dopad na vnímanou latenci systému uživatelem.

\subsubsection*{Kontrola plagiátorství}
\label{subsubsec:kelvin-plagiarism-check}

Vedle evaluační pipeline systém Kelvin provádí rovněž kontrolu plagiátorství.
Ta je ze své podstaty vázána na celou skupinu odevzdání, nikoli na jedno konkrétní, a je proto spouštěna buď manuálně učitelem, nebo automaticky po vypršení termínu odevzdání dané úlohy.
Po spuštění systém porovná každé odevzdání s každým ostatním ve skupině za využití nástrojů DOLOS~\cite{dolos} a MOSS~\cite{moss}.
Díky tomuto uspořádání kontrola plagiátorství nijak neovlivňuje dobu čekání studentů na výsledky automatických testů.

\endinput
